\songtitle{In That Moment}

\tag*{2026}
\begin{notes}
Lyrics written in collaboration between Carlos Thompson and Claude, adapted fron \emph{For That Moment}.
\end{notes}
\begin{style}
trance, house, deep male bass-baritone vocal, half-spoken delivery, vocal subdued in the mix, wall of sound, dense atmospheric synth layers, reverb-heavy pads, pulsing sub-bass, mid-tempo build and drop structure, sensual, intimate, immersive
\end{style}
\begin{lyrics}
\songpart{Intro — sparse, filtered pad, vocal fragment looping, heavily reverbed}
(sung, distant, processed)\\
In that moment... she is all...\\
In that moment... she is all...

\songpart{Breakdown 1 — spoken, flat affect, minimal beat underneath}
Someone I like. A friend from real life.\\
A face I follow. Someone I've never met.\\
Online communities. Online celebs.\\
Doesn't matter where. Doesn't matter who.

\annotation{Rise 1 — sung/chanted, layers building, rhythmic fragment repeating with rising filter}\\
Doesn't matter where she's from...\\
doesn't matter who she is...\\
doesn't matter where... doesn't matter who...

\annotation{Drop 1 — full beat, vocal chop riding the rhythm}\\
She is her.\\
She is her.\\
Only her. Only her.\\
She is her.

\songpart{Breakdown 2 — spoken, pulled back, more space}
What is that feeling?\\
For each one, it's its own.

\annotation{Rise 2 — layered whispers/fragments stacking, building tension}\\
Platonic...\\
or admiration...\\
a lifetime...\\
or a night...

\annotation{Drop 2 — biggest drop, fuller arrangement, hook layered with Drop 1's phrase}\\
For that moment...\\
she is everything.\\
For that moment...
(she is her... she is her...)\\
No before, no after,\\
just for that moment.

\songpart{Outro — stripping back to intro's loop, sung, distant again, processed, fading out}
In that moment... she is all...\\
In that moment... she is all...\\
In that moment... she is all...\\
In that moment... she is all...
...
\end{lyrics}

